\section{Quantitative conventions and symbols}\label{prelim:section}

The analytic argument uses polynomial niltests both to detect structured
correlation and to compare positive functions. This section fixes their
quantitative meaning and develops the operations needed to reduce their
degree. The rational calculus keeps its costs polynomial in logarithmic
complexity when dimensions grow; the symbol calculus separates normalized
real control from rational denominators. We then prove the local box
versions of the inverse and converse inverse theorems and the normalized
square construction used in the step-drop argument.

For a positive integer $M$, write $[M]=\{1,\ldots,M\}$.
If $Q$ is a nonempty finite set, $\E_{x\in Q}$ denotes its
uniform probability average.  All correlations on such a set are
normalized: ``$f$ correlates with $T$ by at least $\rho$'' means
\[
 |\langle f,T\rangle_Q|\geq\rho,\qquad
 \langle f,T\rangle_Q=\E_{x\in Q}f(x)\overline{T(x)}.
\]
Real nonnegative tests require no conjugation.  Slices and subboxes
always carry their own uniform probability measures unless another
measure is specified.

Throughout this section the degree $s$ is fixed.  An expression
$\mathcal P(p)$ denotes $(2+p)^C$, and $\mathcal B(p)$ denotes
$\exp((2+p)^C)$, for a constant $C$ depending only on the fixed degrees.
The constant can increase from one occurrence to the next.  These are
existence bounds.  In particular, ``$T\geq\mathcal B(p)$ sufficiently
large'' means that some bound of the displayed form suffices.  A
polynomial number of products, or a bounded number of compositions of
these bounds, is permitted.  An unbounded number of successive
polynomial changes of the budget is not implicit in this notation.

For $T=(T_1,\ldots,T_n)$ and a multi-index $\alpha$, put
$T^\alpha=\prod_iT_i^{\alpha_i}$.  Every coefficient estimate below uses
these actual side lengths.  No bound on $\max T_i/\min T_i$ is assumed.
The height of a rational array bounds its numerators and positive
denominators.  A \emph{denominator bound} does not bound numerators.
The letter $n$ counts integer parameters.  In the nilmanifold
discussion $d$ denotes the group dimension; later uses of $d$ as a
propagation degree, polynomial weight, or patch rank are specified
locally.  Variable dimensions and ranks enter the displayed budgets,
whereas constants may depend on fixed degrees.

\subsection{Filtered nilmanifolds}

\begin{definition}[Niltests]\label{prelim:niltest}
A rational filtered nilmanifold consists of a connected, simply
connected nilpotent Lie group $G$, a cocompact lattice $\Gamma$, and
connected rational subgroups
\[
 G=G_1\supseteq G_2\supseteq\cdots\supseteq G_{s+1}=\{1\},
 \qquad [G_i,G_j]\subseteq G_{i+j}.
\]
Put $d=\dim G$.  We fix rational exponential coordinates on its Lie algebra
$\mathfrak g$, with rational bases of the filtered subspaces
$\mathfrak g_i=\log G_i$.  A polynomial map $g:\Z^n\to G$ is adapted
to this filtration when $\log g$ modulo $\mathfrak g_{i+1}$ has ordinary
degree at most $i$.  With explicitly assigned variable weights, ordinary
degree is replaced by weighted degree.

A degree-$s$ niltest is $t\mapsto F(g(t)\Gamma)$, where $F$ is bounded
and Lipschitz for the quotient of the right-invariant Riemannian
metric that makes the declared coordinate basis orthonormal at the
identity.  A different metric is permitted only with a comparison
constant whose logarithm is included in the complexity budget.
Its complexity is at most $p$ if the dimension, logarithmic heights of
the structural constants and filtered bases, and
$\log(2+\|F\|_\infty+\Lip F)$ are at most $p$, and if
\begin{equation}\label{prelim:grids}
 l\Z^d\subseteq\log\Gamma\subseteq l^{-1}\Z^d
 \quad\text{for some integer }1\leq l\leq e^p.
\end{equation}
There is no coefficient bound on $g$.  On $\Z/N\Z$ the evaluation
uses representatives $0,\ldots,N-1$.  Degree-zero tests are constants.
The group of dimension zero is allowed.
\end{definition}

Here and below one can enlarge the budget by $\mathcal P(p)$ when
changing rational coordinates.  The following elementary facts explain
the uniformity in a variable dimension.

\begin{lemma}[Rational operations]\label{prelim:rational-operations}
Suppose a polynomial number of rational matrices have polynomially many
rows and columns and entries of height $\mathcal B(p)$.  Their kernels,
images, intersections, quotient maps and sections of surjections have
rational bases or matrices of height $\mathcal B(p)$, provided these
operations are performed jointly, or at a bounded depth.

The analogous assertions hold for rational filtered subgroups,
quotients by connected rational normal subgroups, and coordinate
formulae obtained from the Baker--Campbell--Hausdorff formula at fixed
step.  The resulting lattices have grid bounds of the form
\eqref{prelim:grids} with $l\leq\mathcal B(p)$.
\end{lemma}

\begin{proof}
Clear all denominators by multiplying them.  The logarithm of this
product is polynomially bounded.  For an $r\times r$ integer matrix
whose entries have magnitude at most $H$, the determinant has magnitude
at most $r!H^r$.  Cramer's rule applied to a nonzero maximal minor gives
the asserted bases and sections.  An intersection is a single kernel
calculation on the combined defining equations; performing a separate
elimination for every subspace is unnecessary.

At fixed step the Baker--Campbell--Hausdorff formula has bounded
commutator length.  In dimension $d$, each resulting coordinate has
only $d^{O_s(1)}$ terms.  The same remains true, with the number of
factors included in the polynomial, for a polynomial number of
multiplications.  Apply the preceding rational estimates to their
coefficients.  Restricting or projecting the grids in
\eqref{prelim:grids}, and clearing the denominators of the chosen
rational bases, proves the subgroup and quotient grid assertions.
\end{proof}

The log-polynomial definition in Definition~\ref{prelim:niltest}
agrees with the Taylor/binomial definition of an adapted polynomial
sequence.  Indeed, modulo each successive layer, subtract the
homogeneous terms using factors
$\exp(X\binom{t}{\alpha})$ with $X\in\mathfrak g_{|\alpha|}$.
The finite Baker--Campbell--Hausdorff formula proves the converse and
closure under products.  This also proves that a rational linear
section respecting every filtered layer lifts polynomial maps.  The
section need not be a Lie algebra homomorphism: lifting in logarithmic
coordinates is sufficient.

We shall use finite covers with the same quantitative convention.
Choose an ordered rational basis adapted to the filtration.  After
scaling every basis vector by a sufficiently divisible integer of size
$\mathcal B(p)$, its integral ordered products form a subgroup of
$\Gamma$; this follows directly from the triangular multiplication
and inverse formulae.  This supplies controlled Mal'cev coordinates
and bounded fundamental representatives.  It also supplies common
congruence covers for a polynomial number of rational requirements.
In particular, if $z$ has bounded denominator but arbitrary numerator,
the formula $\Ad_z=\exp(\operatorname{ad}\log z)$ shows that a fixed
sufficiently deep integral grid is contained in $\log(z\Gamma z^{-1})$.
The depth depends on denominators, not numerators.  A bounded-degree
rational polynomial map with common denominator $\mathcal B(p)$ is
periodic modulo such a cover with period $\mathcal B(p)$: subtract its
values at arguments congruent modulo a sufficiently divisible period
and use the same polynomial formulae.

When a cover normal in $\Gamma$ is needed, one can obtain it without
the factorial loss from taking a permutation-group core.  Let
$\Lambda_0\leq\Gamma$ be the controlled cover already constructed,
and choose $m_0\leq\mathcal B(p)$ with
$m_0\Z^d\subseteq\log\Lambda_0$.  Choose an integer $m$ divisible
by $lm_0$, and let
\[
 \Lambda=\langle\gamma^m:\gamma\in\Gamma\rangle.
\]
This subgroup is characteristic in $\Gamma$.  Since
$\log(\gamma^m)=m\log\gamma\in m_0\Z^d$, every generator
belongs to $\Lambda_0$, and hence $\Lambda\subseteq\Lambda_0$.
A Mal'cev normal form with $d$ integral exponents shows that
$[\Gamma:\Lambda]\leq m^d$: in the quotient each exponent can
be reduced modulo $m$, using normality.  Moreover
\[
 ml\Z^d\subseteq\log\Lambda
     \subseteq\log\Gamma\subseteq l^{-1}\Z^d.
\]
Thus normality, index and grid bounds all have the same quantitative
cost.  This construction does not assert that an arbitrary uniform
exponential grid is itself invariant under conjugation.

On bounded fundamental representatives these coordinate formulae
also give Lipschitz bounds $\mathcal B(p)$ for bounded rational group
homomorphisms and bounded left translations.  Fixed right translation
is an isometry before passing to a suitable cover.  For varying left
translations, right invariance gives the useful exact estimate
\[
 d_G(ag,a'g)=d_G(a,a').
\]
Comparing small coordinate balls with metric balls, and constructing
Lipschitz partitions at precision $\mathcal B(p)^{-1}$, therefore costs
only $\mathcal B(p)$.

\begin{lemma}[Reconstruction from rational images]
\label{prelim:reconstruction}
Let $\phi:H\to L$ be a rational homomorphism between controlled
filtered nilpotent groups, and suppose the matrix of its differential
in the declared rational coordinates has height $\mathcal B(p)$.
Assume its kernel is connected and rational.  A controlled
Lipschitz function on $H/\Gamma_H$ invariant under left translation by
$\ker\phi$ can be represented, on $\phi(H)$, by a controlled
Lipschitz function of its image in a sufficiently fine lattice cover
of $L$.  The cost is $\mathcal B(p)$.  The same assertion applies to
a polynomial number of homomorphisms by taking their product.
\end{lemma}

\begin{proof}
Write $J=\phi(H)$.  Rational linear algebra and the grid bounds allow
a target lattice $\Gamma'_L$ with
$J\cap\Gamma'_L\subseteq\phi(\Gamma_H)$.  Thus the function is
well defined on $J/(J\cap\Gamma'_L)$: changing a source lift changes
it by the connected kernel, and changing an image representative
changes it by the image of a source lattice element.

For the Lipschitz estimate take bounded representatives in $J$.
A target lattice element identifying two sufficiently close image
points is itself bounded.  In logarithmic coordinates its distance
from the rational subspace $\log J$ is either zero or at least
$\mathcal B(p)^{-1}$.  This follows by applying a bounded rational
matrix defining that subspace to the bounded-denominator lattice
coordinates.  At a smaller distance it consequently belongs to $J$.
The small displacement then lifts by a bounded rational linear section
of $\log H\to\log J$, which proves the local Lipschitz bound.
Boundedness of the function gives the global bound for larger
distances.  Real and imaginary parts can be extended separately from
the image by the Lipschitz extension theorem; clipping preserves a
specified real interval when the function is real valued.
\end{proof}

\begin{lemma}[Vertical decomposition]\label{prelim:vertical}
Put $e(x)=\exp(2\pi ix)$.  A controlled Lipschitz function can be
approximated uniformly, to error $\mathcal B(p)^{-1}$ with any prescribed
polynomial logarithmic accuracy, by a sum of $\mathcal B(p)$ functions
$F_\eta$ satisfying
\[
 F_\eta(zx)=e(\eta(\log z))F_\eta(x),\qquad z\in G_s.
\]
The frequencies and the Lipschitz complexities have size
$\mathcal B(p)$.  If the initial function has a specified character on
a subtorus, every nonzero summand has that character there.
\end{lemma}

\begin{proof}
The quotient $G_s/(G_s\cap\Gamma)$ is a central torus with a rational
lattice basis of controlled height.  Convolve in that torus with a
product of sufficiently high-order Fej\'er kernels, then expand the
resulting trigonometric polynomial.  If the torus has dimension
$d\leq\mathcal P(p)$, Lipschitz constant $L\leq\mathcal B(p)$,
and required error $\delta\geq\mathcal B(p)^{-1}$, the cutoff in each
coordinate can be taken of size $(1+dL/\delta)^{O(1)}$.  The number of
terms and all frequency heights are thus $\mathcal B(p)$.  Translation
by the central torus commutes with the relevant metric comparisons,
so the summands retain the asserted Lipschitz bounds.  The last
assertion follows by taking the Fourier transform along the subtorus.
\end{proof}

\subsection{Symbols and separation}

Let $\mathfrak n=\bigoplus_{j=1}^s\mathfrak g_j/\mathfrak g_{j+1}$
be the associated graded Lie algebra.  The \emph{symbol} of $g$ is the
map $X$ into its simply connected group for which $(\log X)_j$ is the
degree-$j$ homogeneous part of $\log g$ projected to grade $j$.
Symbol extraction respects products, is unchanged by parameter
translation, and commutes with integer dilation.  These statements
follow by retaining the terms whose parameter degree equals their
Lie layer in the Baker--Campbell--Hausdorff formula.

A symbol is \emph{slow} if its coefficient of $t^\alpha$ has size at
most $\mathcal B(p)T^{-\alpha}$.  A symbol is \emph{rational} if all
its logarithmic coefficients have a common denominator at most
$\mathcal B(p)$; their numerators are unrestricted.  Both classes
are closed under products and inverses at bounded depth.

We distinguish this from a \emph{fully slow} polynomial map on the
box.  In coordinates translated so that the box starts at the
origin, a scalar or vector polynomial $q(t)=\sum_\alpha c_\alpha
t^\alpha$ is fully slow at budget $p$ if
\[
 \|c_\alpha\|\leq\mathcal B(p)T^{-\alpha}
 \quad\text{for every coefficient, including }\alpha=0.
\]
A group-valued polynomial $a$ is fully slow when its polynomial
$\log a$ satisfies these bounds in the declared logarithmic
coordinates.  Thus its constant logarithm is bounded by
$\mathcal B(p)$, and the condition controls every lower-degree
coefficient as well as the homogeneous symbol.  For a translated
box the definition is applied after translating the parameter
coordinates; it always uses that box's actual side lengths.

\begin{lemma}[Separation]\label{prelim:separation}
Let $U$ be a rational subspace of controlled height.  Suppose $a+b\in U$,
where $|a|\leq\mathcal B(p)T^{-\alpha}$, $|\alpha|>0$, and $b$ has
bounded denominator.  If $\min T_i\geq\mathcal B(p)$ is sufficiently
large, then $a,b\in U$.
\end{lemma}

\begin{proof}
Choose a rational matrix $Q$ of controlled height with kernel $U$.
The vector $Qb=-Qa$ has bounded denominator.  A nonzero coordinate
therefore has absolute value at least $\mathcal B(p)^{-1}$, whereas
$|Qa|\leq\mathcal B(p)/\min T_i$.  Taking the lower length cutoff
sufficiently large forces $Qa=Qb=0$.
\end{proof}

\begin{lemma}[Compatible symbol splittings]
\label{prelim:symbol-splitting}
Suppose
\[
 X=EPR=E'P'R',
\]
where $E,E'$ are slow, $R,R'$ are rational, and $P,P'$ take values
in the same graded rational Lie subalgebra $U$ of controlled height,
with its logarithmic height included in $p$.  Then $E^{-1}E'$
and $R'R^{-1}$ take values in $U$, for sufficiently large sides.

If a polynomial number of controlled graded rational Lie subalgebras
$U^{(b)}$ each admit a splitting of $X$ of this kind, then their
intersection admits one, with the same quantitative convention.
\end{lemma}

\begin{proof}
For the first assertion put $A=E^{-1}E'$ and $D=R'R^{-1}$, so that
$P=AP'D$.  Work upwards in grade.  Once the lower-grade logarithmic
coefficients of $A,D$ belong to $U$, every lower-grade commutator in
this identity belongs to $U$.  Its current-grade coefficients therefore
say that a slow coefficient plus a rational coefficient belongs to
$U$.  Lemma~\ref{prelim:separation} completes the induction.

For the second assertion construct common slow and rational factors
one grade at a time.  At a given grade, reset the lower grades of each
individual splitting to the already chosen common factors.  The first
assertion, applied in the truncated group, puts the required ratios
inside $U^{(b)}$.  Lifting those ratios by setting the unused logarithmic
layers to zero lets one absorb them into the middle factor.

Let $V$ be the current coefficient space, and let
$q:V\to\bigoplus_b V/U^{(b)}$ be the simultaneous quotient map.
The remaining coefficient array $v$ satisfies $q(v)=a+b$, with
$a$ slow and $b$ rational.  Modulo the rational subspace $q(V)$,
Lemma~\ref{prelim:separation} puts $a$ and $b$ separately in $q(V)$.
A bounded rational section of $q$ gives common slow and rational
corrections; their removal puts the current coefficients in
$\ker q=\bigcap_bU^{(b)}$.  All matrices are solved jointly at each
grade, and there are only $s$ grades.
\end{proof}

These symbolic splittings lift to actual polynomial maps.  Lift the
homogeneous logarithmic coefficients of the slow and rational symbols
through fixed rational sections, assigning no other logarithmic terms;
call the lifts $e_0,r_0$.  Thus $e_0(0)=r_0(0)=1$.  Choose
$\kappa\in\Gamma$ so that $c=g(0)\kappa^{-1}$ is a bounded
representative, and put
\[
 E=e_0c,\qquad
 b=c^{-1}e_0^{-1}g\kappa^{-1}r_0^{-1},\qquad R=r_0.
\]
Then $g=EbR\kappa$ exactly and $b(0)=1$.  Multiplication by
constants does not change a symbol, so the symbol of $b$ is the
middle symbol of the original splitting.  Fixed-step BCH shows that
$E$ is fully slow, including its constant coefficient, while $R$
has a bounded common denominator.  All these assertions use the
coordinates translated to the box origin.  If the middle symbol takes
values in $U$, then $b$ is adapted to the refiltration
\[
 \mathfrak h_j=\{v\in\mathfrak g_j:
                  v\bmod\mathfrak g_{j+1}\in U_j\}.
\]
Commutator closure follows from that of $U$, and the middle map takes
values in $H_1=\exp\mathfrak h_1$.  This statement does not bound
the lower-degree coefficients of that middle map.

\subsection{Gowers norms and the normalized square}

Here is the precise normalization of the Gowers norms.  Write
$\mathcal C z=\overline z$.  On a finite abelian group $H$,
\[
 \|f\|_{U^j(H)}^{2^j}
 =\E_{x,h_1,\ldots,h_j\in H}
   \prod_{\omega\in\{0,1\}^j}
   \mathcal C^{|\omega|}f(x+\omega\cdot h).
\]
For a finite nonempty integer box $Q\subseteq\Z^n$, put
\[
 C_j(Q)=\#\{(x,h_1,\ldots,h_j)\in(\Z^n)^{j+1}:
                    x+\omega\cdot h\in Q\text{ for every }\omega\}.
\]
Writing $f_0$ for the zero extension of $f$ to $\Z^n$, we fix
\begin{equation}\label{prelim:box-gowers}
 \|f\|_{U^j(Q)}^{2^j}
 =\frac{1}{C_j(Q)}
   \sum_{x,h_1,\ldots,h_j\in\Z^n}
   \prod_{\omega\in\{0,1\}^j}
      \mathcal C^{|\omega|}f_0(x+\omega\cdot h).
\end{equation}
Thus $\|1\|_{U^j(Q)}=1$ and $\|f\|_{U^1(Q)}=|\E_Qf|$.
For fixed $j$ there are constants $0<c_j<C_j'$ such that
\[
 c_j^n|Q|^{j+1}\leq C_j(Q)\leq (C_j')^n|Q|^{j+1}.
\]
The upper bound follows by choosing $x$ and then bounded coordinate
differences; for the lower bound restrict $x$ to a central subbox
and each $h_i$ to a sufficiently small box.  Bounded side lengths
are covered by taking all corresponding increments zero.  Therefore
normalizing the same sum by $|Q|^{j+1}$ changes the norm by at
most $\exp(O_j(n))$.

Embed $Q$, after translation, in
$H=\prod_{i=1}^n\Z/(4T_i)\Z$, where $T_i$ are its side lengths.
All cube identities with vertices in $Q$ are ordinary integer
identities: a face discrepancy cannot be a nonzero multiple of the
corresponding period.  If $\widetilde f$ is the zero extension in
$H$, the exact comparison is
\begin{equation}\label{prelim:ambient-gowers}
 \|\widetilde f\|_{U^j(H)}^{2^j}
 =\frac{C_j(Q)}{|H|^{j+1}}\|f\|_{U^j(Q)}^{2^j}.
\end{equation}
In particular the triangle and mixed Cauchy--Schwarz inequalities
transfer with this fixed normalization.

For clarity, define $\Delta_hf(x)=f(x)\overline{f(x+h)}$.
The derivative identity on a group uses the uniform average
\[
 \|f\|_{U^{j+1}(H)}^{2^{j+1}}
 =\E_{h\in H}\|\Delta_hf\|_{U^j(H)}^{2^j}.
\]
On a box let $Q_h=Q\cap(Q-h)$.  The corresponding exact identity
has cube-count weights:
\begin{equation}\label{prelim:box-derivative}
 C_{j+1}(Q)\|f\|_{U^{j+1}(Q)}^{2^{j+1}}
 =\sum_{h:Q_h\ne\varnothing}
     C_j(Q_h)\|\Delta_hf\|_{U^j(Q_h)}^{2^j}.
\end{equation}
This is obtained by fixing the last increment in the unnormalized
cube sum; in particular $\sum_hC_j(Q_h)=C_{j+1}(Q)$.
We use either the ambient group identity or this weighted box
identity, never an unweighted average of separately normalized
overlap norms.

\begin{theorem}[Quasipolynomial inverse theorem]
\label{prelim:inverse}
Fix an integer $s\ge1$. Let $p\ge2$. A unit-bounded function on an interval
whose $U^{s+1}$ norm is at least $e^{-p}$ correlates by at least
$\mathcal B(p)^{-1}$ with a unit-bounded degree-$s$ niltest of
complexity $\mathcal P(p)$. The same conclusion holds on integer
boxes and on products of cyclic groups, with the following dimension
convention: for every fixed
$A\ge1$, dimension $n\le(2+p)^A$ is permitted, and the constants
and polynomial exponents may depend only on $s,A$. On a product
of cyclic groups the niltest is evaluated in the coordinatewise
integer representatives.
\end{theorem}

The interval assertion is the inverse theorem of Leng, Sah and
Sawhney~\cite[Theorem~1.2]{LSSInverse2024}; for $s=1$ it follows
from Fourier inversion and Parseval. The cited Lipschitz norm
includes the supremum norm. Dividing its correlator by
$\max(1,\|F\|_\infty)$ therefore makes the correlator unit-bounded
at an admissible cost.

We first compare the coordinate conventions. In the same adapted
Mal'cev basis of dimension $d$, let $d_c$ be the source's
right-invariant chain metric and let $d_R$ be our right-invariant
Riemannian metric. For $d\ge1$,
\[
 d_c(x,y)\le d_R(x,y)\le d\,d_c(x,y).
\]
For the second inequality, write $xy^{-1}$ as an ordered product
of coordinate exponentials. Traversing its factors from right to
left gives a path of length at most the sum of the absolute
coordinate values. Apply this to every link in a chain. For the
first inequality, subdivide a smooth path into small displacements;
the first-order coordinate sup norm is bounded by its Riemannian
speed. Take successively the mesh limit and the infimum over paths.
Both comparisons pass to quotient metrics. Fixed-step BCH gives
polynomially many rational coefficients for the changes between
ordered and exponential coordinates. Clearing their denominators
in both directions supplies the lattice grids of
Definition~\ref{prelim:niltest} with polynomial logarithmic cost.
No bound on the polynomial orbit coefficients is required.

Here are the local extensions of the interval theorem. Put
$j=s+1$, translate an integer box to
$Q=\prod_{i=1}^n\{0,\ldots,T_i-1\}$, and set
\[
 b_i=4T_i,\qquad a_1=1,\qquad
 a_i=\prod_{r<i}b_r,\qquad L=\prod_i b_i,
 \qquad \Phi(x)=1+\sum_i a_ix_i.
\]
Then $\Phi(Q)\subseteq[L]$ and $|Q|/L=4^{-n}$.
If $\Phi(u)+\Phi(v)=\Phi(w)+\Phi(z)$, the coordinate discrepancies
$c_i=u_i+v_i-w_i-z_i$ satisfy $|c_i|<b_i$ and
$\sum_i a_ic_i=0$. Reduction modulo $b_1$, followed by division
and repetition, gives $c_i=0$ for every $i$. Thus this affine
map preserves and reflects every two-face identity, and hence
all additive cube identities. If $\widetilde f$ equals $f$ on
$\Phi(Q)$ and is zero elsewhere in $[L]$, the unnormalized cube
sums agree exactly. Consequently
\[
 \|\widetilde f\|_{U^j([L])}^{2^j}
 =\frac{C_j(Q)}{C_j([L])}\|f\|_{U^j(Q)}^{2^j}.
\]
The ratio is at least $\exp(-O_j(n))$ by the cube-count bounds.
The interval input therefore has norm at least
$\exp(-p-O_j(n))$. Pull its correlating orbit back along $\Phi$.
The correlation on $Q$ is the interval correlation multiplied by
$L/|Q|$, and the degree is unchanged. This proves the box assertion,
including anisotropic boxes.

Finally let $H=\prod_{i=1}^n\mathbb Z/N_i\mathbb Z$.
Partition each representative interval into
$\min(4,N_i)$ balanced nonempty intervals. Each piece has diameter
less than $N_i/3$ and length at least $N_i/8$. There are at most
$4^n$ product boxes. The Gowers triangle inequality selects one,
call it $Q$, such that
\[
 \|f1_Q\|_{U^j(H)}\ge4^{-n}\|f\|_{U^j(H)}.
\]
A two-face discrepancy with all vertices in this box has absolute
value less than $2N_i/3$ in coordinate $i$. If it is zero modulo
$N_i$, it is therefore zero as an integer. The group and integer
cube sums agree, so
\[
 \|f1_Q\|_{U^j(H)}^{2^j}
 =\frac{C_j(Q)}{|H|^{j+1}}\|f|_Q\|_{U^j(Q)}^{2^j}.
\]
In particular the restricted box norm is at least $e^{-p}4^{-n}$.
The box result gives a cap-one representative-coordinate niltest
$T$ with $|\E_Q f\overline T|\ge\eta$, where
$\eta\ge\exp(-\operatorname{poly}_s(p+n))$. Put
$\lambda=(|Q|/|H|)\eta\ge8^{-n}\eta$.

We may insert a smooth selector without a lower bound on the
cyclic moduli. For each nonsingleton cyclic coordinate, take the
circle arc whose endpoints lie halfway between the grid points
just inside and outside its selected interval. A $[0,1]$-valued
Lipschitz approximation $W_i$ with transition radius $\theta$
and Lipschitz constant $O(\theta^{-1})$ has discrete $L^1$ error
$O(\theta)$. Indeed, if $\theta<1/(2N_i)$ it agrees with the
indicator on the grid; otherwise each transition interval
contains $O(N_i\theta)$ grid points. For $N_i=1$ take $W_i=1$.
The product selector $W=\prod_iW_i$ has error $O(n\theta)$.
Choose $\theta$ so that this is at most $\lambda/4$.
Then
\[
 |\E_H f\overline{WT}|\ge3\lambda/4.
\]
The selector is a degree-one niltest on an $n$-dimensional torus,
and its logarithmic Lipschitz cost is polynomial in $p+n$.
Thus $WT$ has degree at most $s$, cap one, and the required
complexity and correlation. This completes the local extension.

For a subgroup $H\le G$, write
$H^\Delta=\{(h,h):h\in H\}$ for its diagonal subgroup in $G\times G$.

The relative-square construction and its diagonal quotient occur in
Green and Tao~\cite[Proposition~7.2 and Lemma~7.4]{GreenTaoOrbits2012}.
We record the normalization and its quantitative coordinate proof here
so that the number of parameters and the group dimension remain explicit.

\begin{lemma}[Normalized square]\label{prelim:square}
Let $F$ be vertical of frequency $\eta$ on $G_s$ and let $g$ be an
adapted polynomial map.  For a fixed integer shift $h$, after a
constant right lattice translation and a bounded constant left
translation of $g(t+h)$, the pair with $g(t)$ is polynomial in
\[
 G^\square_j=G_j^\Delta(G_{j+1}\times G_{j+1}).
\]
The function given by the product of the translated $F$ and
$\overline F$ descends by $G_s^\Delta$.  It is consequently a
degree-$(s-1)$ niltest, of controlled complexity independent of the
coefficients of $g$ and of $h$.
\end{lemma}

\begin{proof}
The horizontal difference $g(t+h)g(t)^{-1}\bmod G_2$ is constant in
$t$.  Reduce it modulo the projected lattice.  Thus choose
$\gamma_h\in\Gamma$ and bounded $\epsilon_h\in G$ whose horizontal
logarithms are $m_h,e_h$, respectively, and whose sum is this
difference.  The first path
$\epsilon_h^{-1}g(t+h)\gamma_h^{-1}$ now has the same horizontal
coordinate as the second path.

Put $d_h(t)=\epsilon_h^{-1}g(t+h)\gamma_h^{-1}g(t)^{-1}$.
The logarithm of $d_h$, in original Lie layer $j$, has parameter
degree at most $j-1$.  In the finite
Baker--Campbell--Hausdorff expansion, a term without a shift
difference or inserted constant cancels, while either operation
lowers parameter degree relative to Lie layer by at least one.
The factorization
\[
 (\epsilon_h^{-1}g(t+h)\gamma_h^{-1},g(t))
     =(d_h(t),1)(g(t),g(t))
\]
therefore proves polynomiality in $G^\square_\bullet$.
Right multiplication by $\gamma_h^{-1}$ leaves the original coset
unchanged.  The new function is $F(\epsilon_hx)\overline{F(y)}$;
its Lipschitz bound depends on the bounded $\epsilon_h$ alone.
Central equivariance cancels on $G_s^\Delta$.  The final filtration
layer is that diagonal group, so quotienting reduces the degree.
All structural and metric bounds follow from the preceding rational
operations; only the coefficients of the new orbit can be large.
\end{proof}

\begin{lemma}[Niltests detect Gowers uniformity]
\label{prelim:dual-gowers}
If a unit-bounded function on a box of dimension $\mathcal P(p)$
correlates by at least $e^{-p}$ with a degree-$s$ niltest of complexity
$\mathcal P(p)$, then its $U^{s+1}$ norm is at least
$\mathcal B(p)^{-1}$.  The assertion also holds on cyclic groups,
including products of a bounded number of translated tests in
representative coordinates.
\end{lemma}

The interval case is the converse inverse theorem of Leng, Sah and
Sawhney~\cite[Lemma~B.5]{LSSInverse2024}.  We give the box argument
to keep its dependence on the number of parameters explicit.

\begin{proof}
Induct on $s$, the case $s=0$ being the definition of $U^1$.
Lemma~\ref{prelim:vertical} selects a top mode $T$, which we may
rescale to be unit-bounded, with
$\rho:=|\langle f,T\rangle_Q|\geq\mathcal B(p)^{-1}$.
Put $a=f\overline T$ and
$c_h=\E_{Q_h}\Delta_ha$.  The exact identity is
\begin{equation}\label{prelim:mean-square}
 \rho^2=\sum_{h:Q_h\ne\varnothing}
               \frac{|Q_h|}{|Q|^2}c_h.
\end{equation}
There are at most $2^n|Q|$ shifts.  With
$\tau=2^{-n-2}\rho^2$, overlaps of volume below $\tau|Q|$
contribute at most $\rho^2/4$ in absolute value.  Discarding also
$|c_h|<\rho^2/4$ leaves a total absolute contribution at least
$\rho^2/2$.  Each weight is at most $1/|Q|$, so at least
$\rho^2|Q|/2$ shifts remain.  Each retained overlap has volume
at least $\tau|Q|$ and hence every side at least $\tau$ times
its original side.

The product of the two mode values is the lower-degree niltest
in Lemma~\ref{prelim:square}.  Induction gives a lower bound
$\mathcal B(p)^{-1}$ for $\|\Delta_hf\|_{U^s(Q_h)}$ on every
retained overlap.  Inserting their count and volume bounds into
\eqref{prelim:box-derivative}, and using the bounds for $C_s(Q_h)$,
proves the assertion.  All dimension factors are $\mathcal B(p)$.

In the cyclic case cut at every wrap point of the bounded number of
translates and subdivide further into intervals of diameter less
than $N/3$.  This gives a bounded number of ordinary intervals.
Select a piece retaining a controlled fraction of the correlation.
Every cyclic cube supported on that piece has ordinary integer
face identities, since a face discrepancy has magnitude less
than $2N/3$ and is a multiple of $N$.  Thus the integer and
cyclic cube sums of its zero extension coincide, with the stated
normalization factors.
Finally a cutoff of an interval cannot convert a sufficiently small
$U^j$ norm, $j\geq2$, into a large one at the present costs: smooth
the cutoff with sufficiently small $\mathcal B(p)^{-1}$ $L^1$
error and approximate it by
a Fourier sum of $\mathcal B(p)$ total coefficient mass.  Use
character invariance and the triangle inequality, and bound the
$U^j$ norm of a unit-bounded error by its $L^1$ norm to the power
$1/2^j$.  This transfers the conclusion back to the uncut cyclic
function.
\end{proof}

\section{Dropping the top frequency of a biased symbol}
\label{prelim:step-drop-section}

A biased top-frequency niltest need not have a slowly varying orbit:
large rational coefficients may remain. We separate these contributions
at the level of the homogeneous symbol, leaving a structured factor
whose top layer is annihilated by the tested frequency.

The antecedents are Leng's efficient equidistribution theorem
\cite[Theorem~4]{LengEquidistribution2023} and the degree-reduction
results of Leng, Sah and Sawhney
\cite[Theorem~5.4 and Corollary~5.5]{LSSInverse2024}. Unequal side
lengths already occur in Leng's theorem. The proof here supplies the
symbol formulation with bounds uniform in a growing number of
parameters, retaining each side length throughout.

\begin{lemma}[Step drop for symbols]\label{prelim:step-drop}
Fix $s\geq1$.  Let $F(g(t)\Gamma)$ be a degree-$s$ niltest of
complexity at most $p$, on a box with $n\leq p$ variables and sides
$T_i$.  Suppose $F$ is vertical of top frequency $\eta$ of height
at most $e^p$, and
\[
 \left|\E_tF(g(t)\Gamma)\right|\geq e^{-p}.
\]
For $\min_iT_i\geq\mathcal B(p)$ sufficiently large, the symbol
admits a factorization $X=EPR$, where $E$ is slow, $R$ is rational,
and $\log P$ takes values in a graded rational Lie subalgebra $W$
of height $\mathcal B(p)$ such that $\eta(W_s)=0$.
\end{lemma}

The proof inducts on the fixed degree rather than the group dimension.
Differentiation gives many biased normalized squares, to which the
lower-degree statement applies. After fixing their rational data, we
obtain slow-plus-rational relations for derivatives of one common
symbol. A lattice argument recovers such relations for each coordinate
derivative; the remaining argument removes their slow and rational
parts grade by grade.

\subsection{Recovering coordinate derivatives}

The recovery step below involves the adjoint action of a symbol whose
coefficients need not be bounded. Its identity diagonal blocks nevertheless
force bounded vertical lattice directions to be determined by their
horizontal coordinates. The following lemma makes this assertion
quantitative at the actual, possibly unequal, side lengths.

\begin{lemma}[Lattice extraction]\label{prelim:lattice}
Let $1\leq n,m\leq\mathcal P(p)$, $0\leq a\leq m$,
$T_i>0$, and $l,R,\delta^{-1}\leq\mathcal B(p)$,
where $l$ is a positive integer, $R\geq1$ and $0<\delta\leq1$.
Put $D_T=\operatorname{diag}(T_1,\ldots,T_n)$, and let $L$ be a
positive diagonal $m\times m$ matrix.  Its first $a$ entries equal
$1$, and every remaining entry is at least $T_*:=\min_iT_i$.
Let $M$ be lower unitriangular, with identity on its first
$a\times a$ block, and let $Y:\R^n\to\R^m$ be any linear map.
Suppose the lattice
\begin{equation}\label{prelim:lattice-definition}
 \Lambda=\{(D_T^{-1}h,L(Yh-Mr)):
             h\in\Z^n,\ r\in l^{-1}\Z^m\}
\end{equation}
contains at least $\delta\prod_iT_i$ vectors of norm at most $R$
with distinct $h$-coordinates.  No size bound is imposed on $M$ or
$Y$.  If $T_*\geq\mathcal B(p)$ is sufficiently large, there are
a rational subspace $K\subseteq\R^a$ of height $\mathcal B(p)$ and
linear maps $S,\mathcal R:K\to\R^m$ with
\begin{equation}\label{prelim:lattice-lifts}
 S=M\mathcal R,\qquad
 S(k)_{\mathrm{hor}}=\mathcal R(k)_{\mathrm{hor}}=k,
 \qquad |LS(k)|\leq\mathcal B(p)|k|.
\end{equation}
On a controlled rational basis of $K$, $\mathcal R$ has a common
denominator $\mathcal B(p)$, with no numerator bound.  Moreover
\begin{equation}\label{prelim:lattice-derivatives}
 Ye_i=s_i+Mr_i+S(k_i),\qquad
 |Ls_i|\leq\mathcal B(p)/T_i,
\end{equation}
where $k_i\in K$ and the $r_i$ have a common denominator
$\mathcal B(p)$.
\end{lemma}

\begin{proof}
The map defining \eqref{prelim:lattice-definition} is invertible on
$\R^{n+m}$, so $\Lambda$ is a lattice.  Let $Z$ be the span of its
vectors of norm at most $R$, and put $\Lambda_Z=\Lambda\cap Z$.
If their $h$-projections lay in a proper subspace, one nonzero real
linear equation would determine one coordinate from the others.
The number of integer possibilities would be at most
\[
 (3R)^{n-1}\frac{\prod_iT_i}{T_*},
\]
contrary to the hypothesis when $T_*$ is sufficiently large.
Thus $Z$ projects onto $\R^n$.

Write $d=\dim Z$.  There are $d$ independent lattice vectors of
norm at most $R$ spanning $Z$.  Their integral span has covering
radius at most $dR/2$, by rounding real coefficients to integers.
The larger lattice $\Lambda_Z$ has no larger covering radius.
Its Voronoi fundamental region is consequently contained in a ball
of radius $dR/2$.  Translate that region by the given bounded
lattice vectors.  The translates are disjoint up to null sets and
are contained in a ball of radius $(1+d/2)R$.  Therefore
\begin{equation}\label{prelim:lattice-covolume}
 \covol(\Lambda_Z)\leq C_*/\prod_iT_i,
 \qquad
 C_*:=\delta^{-1}(C(n+m)R)^{C(n+m)},
\end{equation}
for an absolute constant $C$.

Let $V=Z\cap(\{0\}\times\R^m)$ and
$\Lambda_V=\Lambda\cap V$.  The projection in unscaled
$h$-coordinates is a full-rank subgroup $H\leq\Z^n$; put
$I=[\Z^n:H]$.  The lattice $\Lambda_V$ has full rank in $V$:
the projection of a lattice basis of $\Lambda_Z$ is an integer
matrix of rank $n$ in the unscaled $h$-coordinates, and its real
kernel is spanned by rational, hence integral, kernel vectors.
Orthogonally to $V$, the space $Z$ is the graph
of a linear map $A:\R^n\to V^\perp\subseteq\R^m$,
whose input is the scaled coordinate $D_T^{-1}h$.
Lift a basis of $D_T^{-1}H$ to $\Lambda_Z$ and subtract its
$V$-components when calculating the exterior product with a basis
of $\Lambda_V$.  This gives the exact identity
\[
 \covol(\Lambda_Z)=
 \covol(\Lambda_V)\frac{I}{\prod_iT_i}
 \sqrt{\det(\id+A^*A)}.
\]
Writing $E_A=\sqrt{\det(\id+A^*A)}$, we obtain
\begin{equation}\label{prelim:lattice-product}
 \covol(\Lambda_V)I E_A\leq C_*.
\end{equation}
Every singular value of $x\mapsto(x,Ax)$ is at least one, so
$|(x,Ax)|\leq E_A|x|$.

Put $k=\dim V$, and write a basis of the vertical lattice, after
changing signs if necessary, as the columns of $LMR_b$, where
$R_b\in(l^{-1}\Z)^{m\times k}$.  Scan the rows of $R_b$ from
top to bottom, retaining a row exactly when it increases the span.
Let $J$ be the resulting $k$ pivot rows.  Every earlier nonpivot row
is a linear combination of earlier pivot rows.  Hence multiplication
by the lower unitriangular matrix $M$ changes the selected rows by
a lower unitriangular row operation.  In particular,
\[
 \det(MR_b)_J=\det(R_b)_J.
\]
If $w_j$ denotes the $j$th diagonal entry of $L$, Cauchy--Binet
therefore gives
\begin{equation}\label{prelim:lattice-pivot}
 \covol(\Lambda_V)\geq
 |\det(LMR_b)_J|
 =\Bigl(\prod_{j\in J}w_j\Bigr)|\det(R_b)_J|
 \geq l^{-k}\prod_{j\in J}w_j.
\end{equation}
If any pivot were nonhorizontal, the last expression would be at
least $l^{-k}T_*$.  This contradicts
\eqref{prelim:lattice-product} once $T_*>C_*l^m$.  All pivots are
horizontal.  Thus horizontal projection is injective on $V$ and
its image is a rational subspace $K\subseteq\R^a$.  Equations
\eqref{prelim:lattice-product}--\eqref{prelim:lattice-pivot} give
\[
 I,E_A\leq C_*l^k,\qquad
 l^{-k}\leq\covol(\Lambda_V)\leq C_*.
\]

Let $H_b$ be the horizontal rows of $R_b$.  They are also the
horizontal rows of $LMR_b$.  Every maximal minor of $H_b$ is
bounded in absolute value by $C_*$, and belongs to $l^{-k}\Z$.
Choose a nonzero maximal minor $(H_b)_{J_0}$.  The columns of
$H_b((H_b)_{J_0})^{-1}$ form a rational basis of $K$.  Its
entries are ratios of maximal minors, hence have controlled height.
No short basis for $\Lambda_V$ is required here.

The unique graph lift from $K$ to $V$ has volume expansion
\[
 \frac{\covol(\Lambda_V)}{\covol(H_b\Z^k)}
 \leq C_*l^k.
\]
As its horizontal part is the identity, its singular values are
at least one and this also bounds its norm.  Denote the lift by
$k_0\mapsto LS(k_0)$, and define
\[
 \mathcal R\bigl(H_b((H_b)_{J_0})^{-1}u\bigr)
 =R_b((H_b)_{J_0})^{-1}u.
\]
Then $S=M\mathcal R$.  The matrix on the right has denominator
dividing the nonzero integer $\det(l(H_b)_{J_0})$, whose magnitude
is at most $C_*l^k$.  This proves \eqref{prelim:lattice-lifts}.

Finally $I\Z^n\subseteq H$.  Lift $Ie_i$ to $\Lambda_Z$ and
divide the lift by $I$.  The resulting vector is
\[
 (T_i^{-1}e_i,L(Ye_i-Mr'_i/I)),\qquad
 r'_i\in l^{-1}\Z^m.
\]
Its component orthogonal to $V$ has norm at most $E_A/T_i$.
Write that component's vertical coordinate as $Ls_i$, and the
$V$-component as $LS(k_i)$.  Taking $r_i=r'_i/I$ proves
\eqref{prelim:lattice-derivatives}.  Its denominator divides $lI$.
Zero-dimensional lattices have covolume one; the same argument,
with empty bases, covers $k=0$.
\end{proof}


\subsection{Proof of the step-drop lemma}


\begin{proof}
If $\eta=0$, take $W=\mathfrak n$ and $E=R=1$.  A zero-variable
symbol is also trivial, so assume that the box has at least one
variable and that $\eta\ne0$.
We induct on $s$.  If $s=1$, the group is a vector group and
verticality gives $F(x)=c\,e(\eta(x))$ with $|c|\leq e^p$.
Choose a rational
vector $v$ of controlled height with $\eta(v)=1$.  Write
$\log X(t)=\sum_i a_it_i$.  The product formula for the mean of
a linear phase and the geometric-sum bound give integers $q_i$
such that
\[
 |\eta(a_i)-q_i|\leq e^{2p}/(2T_i).
\]
Indeed the product of the $n$ normalized geometric sums has
magnitude at least $e^{-2p}$, so each individual magnitude is
at least $e^{-2p}$.  Take slow coefficients
$(\eta(a_i)-q_i)v$, rational coefficients $q_iv$, and middle
coefficients in $\ker\eta$.  This proves the base case.

\paragraph{The differentiated symbol.}
Assume $s\geq2$.  Apply the weighted-shift selection in the proof
of Lemma~\ref{prelim:dual-gowers}, after rescaling $F$ to be
unit-bounded.  Together with Lemma~\ref{prelim:square}, it gives at least
$\mathcal B(p)^{-1}\prod_iT_i$ distinct shifts $h$, with
$|h_i|\leq T_i$, for which the normalized square test has a mean
of magnitude at least $\mathcal B(p)^{-1}$ on its overlap, and whose
individual side lengths are at least $\mathcal B(p)^{-1}T_i$.

Write $\mathfrak n_j=\mathfrak g_j/\mathfrak g_{j+1}$.
The associated graded Lie algebra of the quotient square has
diagonal layers $\mathfrak n_j$, $1\leq j<s$, and relative
layers $\mathfrak n_{j+1}$ in square degree $j$.  The relative
part is abelian in this associated graded algebra: the bracket
of two relative inputs has one extra filtration degree.  The
diagonal acts on it by the ordinary adjoint action.

Use the normalization in Lemma~\ref{prelim:square}, writing
\[
 (\log X)_1(h)=e_h+m_h,
\]
where $e_h$ has controlled size and $m_h$ lies in a fixed rational
grid.  For a symbol $C$ write
$Y_h(C)=(\partial_hC)C^{-1}$ for its right logarithmic derivative.
In relative-then-diagonal semidirect coordinates, the square symbol is
\begin{equation}\label{prelim:square-symbol}
 \left(X_{\leq s-1},\,
       Y_h(X)-e_h-\Ad_Xm_h\right).
\end{equation}
Here the relative component in original layer $j$ is its
homogeneous parameter part of degree $j-1$, and the horizontal
part is zero.  Here is a formal verification, including the case
where the original filtration has no grading.  Choose complements
$E_j$ of $\mathfrak g_{j+1}$ in $\mathfrak g_j$ and let
$\delta_z$ act by $z^j$ on $E_j$.  The brackets
\[
 [u,v]_z=\delta_z[\delta_z^{-1}u,\delta_z^{-1}v]
\]
are polynomial in $z$: a component in $E_k$ of a bracket of
$E_i,E_j$ is multiplied by $z^{k-i-j}$, with $k\geq i+j$.
At $z=0$ this is the associated graded bracket.  Write
$\exp_z,\log_z$ for the finite Baker--Campbell--Hausdorff group
with this bracket.  The polynomial
\[
 P_z(t)=\delta_z\log g(t/z)
\]
has constant term $P_0(t)=\log X(t)$.  Similarly the scaled
normalizing constants are $ze_h+O(z^2)$ and $zm_h+O(z^2)$.
Here $O(z^2)$ denotes terms divisible by $z^2$ in a formal
polynomial, without any bound on their coefficients.
The scaled relative product is therefore
\[
 \exp_z(-ze_h+O(z^2))\exp_z(P_z(t+zh))
 \exp_z(-zm_h+O(z^2))\exp_z(-P_z(t)).
\]
When the shift and inserted constants are zero, the middle
product equals the identity for every $z$.  Consequently the
first-order variations of the bracket itself and of the
subleading coefficients in $P_z$ cancel identically.  The
coefficient of $z$ depends only on the displayed shift and
insertions, evaluated in the graded group at $z=0$.
It can thus be calculated over the dual numbers
$\R[\varepsilon]/(\varepsilon^2)$, where the exact identity is
\[
 \exp(-\varepsilon e_h)X(t+\varepsilon h)
 \exp(-\varepsilon m_h)X(t)^{-1}
 =\exp\!\left(\varepsilon
       (Y_h(X)-e_h-\Ad_Xm_h)\right).
\]
Indeed $X(t+\varepsilon h)X(t)^{-1}=\exp(\varepsilon Y_h(X))$,
conjugation of the right insertion gives $\Ad_Xm_h$, and
brackets of two $\varepsilon$ terms vanish.  In original layer
$j$, the coefficient of $z$ is exactly the homogeneous
parameter part of degree $j-1$ of the relative logarithm, since
$\delta_z$ multiplies that monomial by $z^{j-(j-1)}$.
This proves
\eqref{prelim:square-symbol}.  The same calculation gives the
semidirect multiplication rule $y+\Ad_Xy'$ under a preceding
diagonal coordinate $X$.

\paragraph{A common fast subgroup.}
Apply Lemma~\ref{prelim:vertical} to the quotient square.  One
top mode retains bias $\mathcal B(p)^{-1}$, and its restriction
to the relative copy of $\mathfrak n_s$ is $\eta$.
The induction hypothesis factors its symbol into slow, fast and
rational factors.  Every resulting fast subalgebra has polynomial
dimension and rational height $\mathcal B(p)$.  There are only
$\mathcal B(p)$ possibilities, since a rational basis is specified
by polynomially many bounded-height entries.  Retain a set of
$\mathcal B(p)^{-1}\prod_iT_i$ shifts with the same fast
subalgebra $\widetilde U$.

Let $U$ be its diagonal projection, with $U_s=\mathfrak n_s$,
and let $V_j\subseteq\mathfrak n_j$ be its kernel in the
relative part, setting $V_1=0$.  Then
\begin{equation}\label{prelim:UV}
 \eta(V_s)=0,\qquad [U_i,V_j]\subseteq V_{i+j}.
\end{equation}
All these are controlled rational spaces.  For each retained
shift, denote the diagonal splitting by $X=A_hB_hD_h$ through
grade $s-1$.  Fix one retained shift and denote its outside
diagonal factors by $A,D$.  Lemma~\ref{prelim:symbol-splitting}
puts $A^{-1}A_h$ and $D_hD^{-1}$ in $U$.  Lift their logarithms
to $\widetilde U$ by a bounded rational section of the diagonal
projection.  The first lift is slow, and the second is rational.
Absorbing them into the middle factor makes its diagonal equal to
\[
 B=A^{-1}XD^{-1}
\]
for every retained shift.  Extend $A,D$ to original grade $s$
with zero logarithmic component there; the same formula defines
the full symbol $B$ and puts it in $U$.

In a subgroup of a semidirect product with abelian kernel, two
elements with the same diagonal have relative coordinates equal
modulo the kernel.  The fast relative coordinate, denoted $Z_B$,
is therefore common modulo $V$.  Comparing the relative
coordinates in \eqref{prelim:square-symbol} and using
\[
 Y_h(ABD)=Y_h(A)+\Ad_A Y_h(B)+\Ad_{AB}Y_h(D)
\]
gives
\begin{equation}\label{prelim:common-derivative}
 Y_h(B)=s_h+\Ad_Br_h+Z_B\pmod V.
\end{equation}
For clarity, if $u_h,Z_h,v_h$ are the three relative factors,
the two sides before solving for $Y_h(B)$ are
\[
 Y_h(X)-e_h-\Ad_Xm_h
       =u_h+\Ad_AZ_h+\Ad_{AB}v_h.
\]
Thus one can take
\[
 \begin{split}
 s_h&=\Ad_{A^{-1}}(e_h+u_h-Y_h(A)),\\
 r_h&=v_h+\Ad_Dm_h-Y_h(D).
 \end{split}
\]
The coefficients of $s_h$ in original layer $j$, parameter degree
$j-1$, satisfy $\mathcal B(p)T^{-\alpha}$ bounds.  This uses
$|h_i|\leq T_i$ when differentiating the slow symbol $A$.
The coefficients of $r_h$ have bounded denominator; integrality
of $h$ is used for $Y_h(D)$.  These claims include the horizontal
constant terms.  Pigeonhole the denominators and subtract one
retained instance of \eqref{prelim:common-derivative} from every
other instance.  The term $Z_B$ disappears, the number of distinct
shift differences is unchanged up to one, and the side and
coefficient bounds only change by controlled factors.

\paragraph{Extraction of a horizontal kernel.}
Let $\mathcal M$ be the direct sum, over $j=1,\ldots,s$, of
the coefficient spaces of degree-$(j-1)$ homogeneous polynomials
with values in $\mathfrak n_j/V_j$.  Choose controlled rational
coordinates.  The operator $M=\Ad_B$ on $\mathcal M$ is block
lower unitriangular by Lie layer, and is well defined by
\eqref{prelim:UV}.  The linear map $h\mapsto Y_h(B)$ takes
values in $\mathcal M$.  Let $L_T$ multiply a coefficient of
$t^\alpha$ by $T^\alpha$.  The horizontal scales are one and
every other scale is at least $\min_iT_i$.

The subtracted identities supply the bounded lattice vectors
required by Lemma~\ref{prelim:lattice}, with a common grid for
the $r_h$; changing from the chosen rational coordinates to a
rectangular grid merely enlarges its denominator.  Consequently
there are a controlled rational $K\subseteq\mathfrak n_1$,
linear maps $S,R$ on $K$, and, for every coordinate derivative,
\begin{equation}\label{prelim:derivative-system}
 \begin{split}
 &S=\Ad_B R\pmod V,\qquad S(k)_1=R(k)_1=k,\\
 &|L_TS(k)|\leq\mathcal B(p)|k|,\\
 &Y_i(B)=s_i+\Ad_Br_i+S(k_i)\pmod V,\qquad
 |L_Ts_i|\leq\mathcal B(p)/T_i.
 \end{split}
\end{equation}
Here $R$ on a controlled rational basis and all the $r_i$ have
bounded common denominator, and $k_i\in K$ need not be rational.
We henceforth use these identities in the quotient coefficient
space, so equalities of components are always interpreted modulo
$V_j$.

\paragraph{The target subalgebra.}
We have obtained the derivative system \eqref{prelim:derivative-system},
with $B$ taking values in $U$. Define
\[
 W_1=\{v\in K\cap U_1:[v,K]\subseteq V_2\},\qquad
 W_j=V_j\cap U_j\quad(2\leq j\leq s).
\]
These controlled rational spaces form a graded Lie subalgebra.
Indeed $[W_1,W_1]\subseteq V_2\cap U_2$, while every other
required bracket lies in the appropriate $U$ layer and in $V$
by \eqref{prelim:UV}. Also $\eta(W_s)=0$.
It remains to remove slow and rational factors from $B$ until its
logarithmic coefficients lie in $W$.

Write $P_j=(\log B)_j$, updating this notation after each removal.
First we put $P_1$ in $U_1\cap K$. At stage $j\ge2$ we arrange
$[P_{j-1},K]\subseteq V_j$ and then put $P_j$ in $U_j\cap V_j$.
The accompanying derivative terms must vanish at the same time so
that the next layer reduces to an ordinary coefficient equation.
We begin with the two operations that preserve the derivative system.

\paragraph{Removing slow and rational derivatives.}
The system \eqref{prelim:derivative-system} is preserved when
$B=aB'b$, where $a$ is slow and $b$ rational, both in $U$.
The correct transformed quantities are
\begin{equation}\label{prelim:transformation}
 \begin{aligned}
 S'&=\Ad_{a^{-1}}S,& R'&=\Ad_b R,\\
 s'_i&=\Ad_{a^{-1}}(s_i-Y_i(a)),&
 r'_i&=\Ad_b r_i-Y_i(b).
 \end{aligned}
\end{equation}
The identities follow from the product rule for $Y_i$.
Their quantitative bounds persist: after coefficient scaling by
$T^\alpha$, multiplication by a slow coefficient only multiplies
norms by a controlled factor; rational products preserve bounded
denominators.  Also, for $k_s,k_r\in K$, the replacement
\begin{equation}\label{prelim:absorption}
 s_i\leftarrow s_i-S(k_s),\qquad
 r_i\leftarrow r_i-R(k_r),\qquad
 k_i\leftarrow k_i+k_s+k_r
\end{equation}
preserves the system.  Its bounds persist if
$|k_s|\leq\mathcal B(p)/T_i$ and $k_r$ has bounded denominator.

The coefficient of $t_i$ in $P_1$ is
$s_{i,1}+r_{i,1}+k_i\in U_1$.
Modulo the rational space $U_1+K$, separation puts $s_{i,1}$
and $r_{i,1}$ separately in that space.  Use a controlled rational
decomposition into a $U_1$ part and a $K$ part.  Absorb the
$K$ parts by \eqref{prelim:absorption}, and remove the $U_1$
parts by factors $a,b$ whose logarithms have only layer one.
We have arranged
\begin{equation}\label{prelim:horizontal-normalization}
 P_1\in U_1\cap K,\qquad s_{i,1}=r_{i,1}=0.
\end{equation}

We now process $j=2,\ldots,s$.  At the start of stage $j$,
suppose the following hold:
\[
 \begin{array}{ll}
 P_d\in U_d\cap V_d &(2\leq d<j),\\
 {}[P_d,K]\subseteq V_{d+1}&(1\leq d\leq j-2),\\
 S_d=R_d=0 &(2\leq d<j),\\
 s_{i,d}=r_{i,d}=0 &(1\leq d<j).
 \end{array}
\]
The last two lines are statements in the quotient by $V$.
Expanding $S=\Ad_B R$ in layer $j$ gives
\begin{equation}\label{prelim:bracket-system}
 [P_{j-1},k]=S_j(k)-R_j(k)\pmod {V_j},
 \qquad k\in K.
\end{equation}
Indeed all lower terms of $R$ except its horizontal part vanish
modulo $V$, and every term in $\exp(\operatorname{ad}P)k$
with an inner bracket from a layer at most $j-2$ belongs to $V$.

Let $E=U_{j-1}\cap V_{j-1}$, or $E=U_1\cap K$ when $j=2$.
On a fixed controlled rational basis $(k_b)$ of $K$, define the
rational linear map
\[
 Q_j:E\longrightarrow\bigoplus_b\mathfrak n_j/V_j,
 \qquad v\longmapsto([v,k_b])_b.
\]
The right side of \eqref{prelim:bracket-system} is a slow
coefficient array minus a rational one.  Project modulo
$Q_j(E)$, and apply Lemma~\ref{prelim:separation} coefficient
by coefficient.  Each part lies separately in the image.
A controlled rational section provides homogeneous polynomials
$a_{j-1},b_{j-1}$ with values in $E$ such that
\[
 Q_j(a_{j-1})=(S_j(k_b))_b,\qquad
 Q_j(b_{j-1})=-(R_j(k_b))_b.
\]
The first is slow and the second rational.  Remove their
exponentials from the left and right of $B$.  Formula
\eqref{prelim:transformation} now gives
\[
 [P_{j-1},K]\subseteq V_j,\qquad S_j=R_j=0.
\]
Lower vanishing conditions persist.  When $j=2$, the operation
creates horizontal parts $-\partial_i a_1$ and
$-\partial_i b_1$ in $s_i,r_i$; both lie in $K$.
Absorb them by \eqref{prelim:absorption}.  Since the new
$S_2,R_2$ are already zero modulo $V_2$, this restores
\eqref{prelim:horizontal-normalization} without changing their
vanishing.  For $j>2$, the analogous new parts lie in $V_{j-1}$
and are already zero in the quotient.

At this point the right logarithmic derivative satisfies
\[
 Y_i(B)_j=\partial_iP_j\pmod {V_j}.
\]
To verify this, expand it as a finite sum of nested brackets of
$P$ and $\partial_iP$.  A derivative in a preceding layer at
least two lies in $V$.  A horizontal derivative lies in $K$,
and its innermost bracket lies in $V$ by the bracket conditions
just established.  Subsequent brackets preserve $V$ by
\eqref{prelim:UV}.  Likewise $(\Ad_Br_i)_j=r_{i,j}$ modulo
$V_j$, since all its lower components are already zero there.
Thus \eqref{prelim:derivative-system} reduces to
\[
 \partial_iP_j=s_{i,j}+r_{i,j}\pmod {V_j}.
\]
For every monomial $t^\beta$ of degree $j$, choose one $i$ with
$\beta_i>0$ and divide the corresponding derivative coefficient
by $\beta_i$.  This expresses its coefficient, modulo $V_j$,
as a slow part of size $\mathcal B(p)T^{-\beta}$ and a
rational part.  Because $P_j$ takes values in $U_j$, separation
modulo $U_j+V_j$ and a bounded rational section allow both parts
to be lifted into $U_j$.  Remove their exponential lifts at
layer $j$.  We have obtained $P_j\in U_j\cap V_j$.
The derivative identity then says $s_{i,j}+r_{i,j}=0$ modulo
$V_j$; separation, now at scale $T_iT^\alpha$, makes the two
terms vanish separately.  This completes stage $j$.

The final logarithmic coefficients now lie in the target algebra
$W$, whose top layer is annihilated by $\eta$. There were only
$s$ stages, with
joint rational linear algebra at each stage.  All corrections
therefore combine into slow and rational outside factors with
the stated bounds.  Together with the original $A,D$, these
give the required factorization of $X$.
\end{proof}
